% Generated by analysis/build-paper.mjs on data/responses (0 valid response file(s)). Do not edit by hand.
\begin{quote}\textbf{No survey responses have been collected (n = 0).} Every estimate below is shown as \textquotedblleft{}--\textquotedblright{} and will be filled in by running \texttt{npm run paper} on exported response files. Nothing in this section is an empirical result.\end{quote}
\begin{table*}[t]
\caption{How many participants held each category of wrong mental model, among those who had a chance to show it (at least one part they did not skip had an item tagged with that category). Cells show participants with the belief over participants with the chance, and the percentage. The All column adds a 95\% Wilson interval. The last column is A opaque minus B transparent, with a Newcombe hybrid-score 95\% interval, in percentage points.}\label{tab:resp-incidence}
\footnotesize
\begin{tabular}{lp{3.0cm}rp{3.3cm}rrp{2.4cm}}
\toprule
Code & Category & Cases & All & A opaque & B transparent & A $-$ B\\
\midrule
WM1 & Uneven language competence read as ``broken'' & 3 & --  & --  & --  & --\\
WM2 & Undiagnosed fault source & 15 & --  & --  & --  & --\\
WM3 & No model of actual capabilities & 4 & --  & --  & --  & --\\
WM4 & Context limits read as forgetting & 4 & --  & --  & --  & --\\
WM5 & Trust miscalibration & 4 & --  & --  & --  & --\\
WM6 & Mislocated intelligence & 16 & --  & --  & --  & --\\
WM7 & Silent system change & 3 & --  & --  & --  & --\\
LS & Assumed learning or adaptation & 10 & --  & --  & --  & --\\
\bottomrule
\end{tabular}
\par\smallskip{\footnotesize Incidence depends on these scenarios and on how the closed items were built, so it is not a population base rate (\S\ref{sec:survey}). People who used both interfaces are left out of the A and B columns.}
\end{table*}


\begin{table*}[t]
\caption{Incidence of each category by self-reported background stratum (participants with the belief over participants with the chance, with 95\% Wilson intervals), plus a Pearson $\chi^2$ test of independence across strata and Cram\'er's $V$ as the effect size.}\label{tab:resp-strata}
\footnotesize
\begin{tabular}{lp{2.6cm}p{2.5cm}p{2.5cm}p{2.5cm}rrr}
\toprule
Code & Category & $G1$ & $G2$ & $G3$ & $\chi^2$ & $p$ & $V$\\
\midrule
WM1 & Uneven language competence read as ``broken'' & -- & -- & -- & -- & -- & --\\
WM2 & Undiagnosed fault source & -- & -- & -- & -- & -- & --\\
WM3 & No model of actual capabilities & -- & -- & -- & -- & -- & --\\
WM4 & Context limits read as forgetting & -- & -- & -- & -- & -- & --\\
WM5 & Trust miscalibration & -- & -- & -- & -- & -- & --\\
WM6 & Mislocated intelligence & -- & -- & -- & -- & -- & --\\
WM7 & Silent system change & -- & -- & -- & -- & -- & --\\
LS & Assumed learning or adaptation & -- & -- & -- & -- & -- & --\\
\bottomrule
\end{tabular}
\par\smallskip{\footnotesize $^\dagger$ An expected cell count is below 5, so the $\chi^2$ approximation is unreliable and that row needs an exact test (\S\ref{sec:reliability}). Read $V$ and the intervals, not only $p$. Score by stratum (mean percent of closed points, $n$): G1 --; G2 --; G3 --.}
\end{table*}

